\documentclass[journal]{IEEEtran}

\usepackage{amsmath,amssymb,amsfonts}
\usepackage{graphicx}
\usepackage{booktabs}
\usepackage{microtype}
\usepackage[hidelinks]{hyperref}

\graphicspath{{figures/}}

\newcommand{\revps}{\ensuremath{\,\mathrm{rev/s}}}
\newcommand{\fsenv}{\ensuremath{f_{\mathrm{env}}}}

\begin{document}

\title{Coupled Vold--Kalman Order Tracking for Blind Rotor-Speed Annotation
  from Multirotor Ego-Noise}

\author{Dmitrii~Mukhutdinov%
  \thanks{D.~Mukhutdinov is with the Centre for Intelligent Sensing, Queen
    Mary University of London, London, UK (e-mail: flyingleafe@gmail.com).}}

\maketitle

\begin{abstract}
Per-rotor rotational speed (RPS) is the single most informative side channel
for drone audition: it parameterizes the harmonic structure of ego-noise and
conditions modern enhancement models. In practice, however, synchronized
telemetry is rarely recorded, and the datasets that dominate the field ship
either no rotor-speed track or one that is too coarse to align. We present a
coupled Vold--Kalman (VK) order-tracking method that recovers all four rotor
trajectories of a quadrotor directly from its acoustic ego-noise, with no
tachometer and no telemetry initialization. The method couples every harmonic
of every rotor in a single least-squares envelope solve, fuses per-harmonic
phase-slope observations into one frequency update per rotor with
Fisher-derived weights, and anneals the harmonic range to widen the capture
region. A comb-structured blind seeding stage --- a whitened, band-limited
matched-filter scan followed by a residual re-scan and an occupancy-based
stage guard --- removes the need for any external initialization. On the
DREGON free-flight recordings the blind tracker attains pooled per-rotor mean
absolute errors of 0.68--0.74\revps; on an unrelated DJI Matrice~100
recording it reaches 1.03\revps{} with all four rotors captured. The same
recordings defeat neural regressors trained on thousands of hours of
simulated and real mixtures by a factor of 2--4. The tracker runs at 0.04
real-time on a laptop CPU core in its telemetry-refinement configuration and
0.35 real-time blind.
\end{abstract}

\begin{IEEEkeywords}
Order tracking, Vold--Kalman filter, drone audition, ego-noise, rotor speed
estimation, harmonic analysis.
\end{IEEEkeywords}

\section{Introduction}

\IEEEPARstart{A}{coustic} sensing from multirotor platforms is dominated by
the platform's own noise. The broadband component of ego-noise is
approachable with conventional enhancement, but the harmonic component ---
combs of tones at integer multiples of each rotor's blade-passage frequency,
sweeping with every control action --- is both the strongest and the most
structured part of the interference. That structure is an opportunity: if
the per-rotor speed trajectories $r_i(t)$ are known, the frequencies of
every harmonic line are known, and both classical projection methods and
rotor-conditioned neural models exploit them to
advantage~\cite{yen2023ncm,gulli2025rotorcond}.

The catch is that $r_i(t)$ is almost never available. Consumer platforms do
not expose synchronized per-rotor telemetry; research datasets either omit it
or record it at rates and alignments that are unusable for phase-coherent
processing. Of the corpora in common use, DREGON~\cite{strauss2018dregon} is
nearly unique in shipping motor telemetry at all, and even there the
commanded-speed track freezes before landing and must be treated with care.
The practical question this paper addresses is therefore: \emph{can accurate
per-rotor speed trajectories be recovered from the acoustic recording alone,
reliably enough to serve as ground-truth annotation for datasets that lack
telemetry?}

We answer in the affirmative with a method built on the Vold--Kalman (VK)
order-tracking filter~\cite{vold1993high}, adapted in three substantial ways
to the multirotor acoustic setting. First, where classical VK practice
extracts one order at a time, or a handful of orders against a common
tachometer reference, our formulation solves for the complex envelopes of
\emph{all} harmonics of \emph{all} rotors jointly, with explicit coupling
terms between spectrally adjacent tracks. A quadrotor in trim flight
operates its four rotors within a few rev/s of one another; their combs
interlace, and any per-track method must either ignore the resulting
explaining-away problem or resolve it. Second, VK filtering presupposes
known carrier frequencies, whereas here the frequencies are the quantity of
interest. We close the loop with an outer iteration that reads the residual
frequency error of each track from the phase slope of its recovered
envelope, fuses these observations across harmonics, rotors and microphone
channels under physically derived weights, and updates a single latent
trajectory per rotor. Third, we initialize blind: a whitened matched-filter
scan over a comb dictionary, restricted to the frequency band where combs
are resolvable, followed by a re-scan of the residual spectrum and a
per-track guard that reverts refinement stages which destroy a captured
track.

The combination is, to our knowledge, the first demonstration that
order-domain analysis of the class used in mechatronic fault diagnosis
transfers to fully blind acoustic annotation of multirotor flight data at
accuracies below 1\revps{} --- and that it does so on unmodified,
publicly available recordings, at CPU speeds compatible with batch
annotation of entire corpora.

\section{Related Work}

\subsection{Vold--Kalman order tracking}

The VK filter was introduced for high-slew-rate order extraction in
automotive testing~\cite{vold1993high} and extended shortly after to
simultaneous tracking of orders referenced to multiple independent
shafts~\cite{vold1997multi}; its filter characteristics were analyzed
in~\cite{herlufsen1999characteristics,tuma2000characteristics}. The
second-generation formulation minimizes a data-fidelity term plus a
difference-operator smoothness penalty on complex envelopes, which is the
form we build on. A line of adaptive variants followed: angular-displacement
and angular-velocity reformulations~\cite{pan2007adaptive,pan2010extended},
kernel corrections for the imaginary-part error of incomplete complex
bases~\cite{zhu2023splitcossin}, and ridge-detection-initialized IF
estimation via synchrosqueezing~\cite{wang2020vkifsep}. Closest to our
refinement loop is the iterative adaptive VK filter (IAVKF) of Jiang et
al.~\cite{jiang2024iavkf}, which alternates envelope extraction with an IF
correction computed from the unwrapped envelope phase and adapts the filter
bandwidth through an orthogonality fixed point. Our outer loop shares the
same core observation --- the recovered envelope's phase slope \emph{is} the
IF error --- but differs structurally: IAVKF extracts independent components
one at a time by greedy peeling, whereas our components are constrained to
integer combs of four latent fundamentals and are solved jointly with
explicit cross-track coupling. The distinction matters precisely in our
regime, where dozens of harmonics from four near-coincident combs overlap
and greedy subtraction is unreliable.

\subsection{Tacholess and multiharmonic frequency tracking}

Recovering the reference speed without a tachometer is a recognized problem
in machine monitoring; see~\cite{peeters2019tacholess} for a review and,
e.g.,~\cite{wang2019tlot} for a mode-decomposition approach to instantaneous
angular speed estimation. These methods, like ours, extract speed from the
signal itself, but they target a single shaft and typically operate on
vibration signals with favorable SNR. Multiharmonic trackers with explicit
state-space structure go back at least to extended Kalman formulations of
comb tracking~\cite{streit1992multiharmonic}; their capture behavior ---
convergence only within a basin around the true fundamental --- motivates
the annealing schedule we use. In the audio domain, fundamental-frequency
estimation is a mature field, and modern neural estimators are robust in
noise~\cite{han2014pitchnn,riou2023pesto}; polyphonic variants recover
multiple concurrent pitches~\cite{bittner2022basicpitch}. We evaluated
representatives of this family as baselines and found them inadequate for
the task: rotor combs are not speech- or music-like, the four fundamentals
are separated by fractions of a semitone, and per-frame salience methods do
not deliver the sub-\revps{} trajectory accuracy that phase-coherent
processing requires.

\subsection{Drone audition}

Ego-noise suppression on multirotors spans blind source
separation~\cite{wang2020bss}, learned time-frequency
masking~\cite{wang2020dnntf}, and a rapidly growing single-channel deep
enhancement literature~\cite{mukhlish2023survey,liu2025edgebsrof}. Several
recent systems consume rotor-speed side information explicitly: noise
covariance estimation driven by rotor speeds~\cite{yen2023ncm} and
rotor-conditioned enhancement networks~\cite{gulli2025rotorcond} both report
consistent gains from the conditioning signal. This strand is the consumer
of our method's output: high-quality $r_i(t)$ annotation is exactly what
those systems assume and what public data mostly lacks.

\section{Problem Setting and Data}
\label{sec:data}

We consider single- or multi-channel recordings of a quadrotor in flight.
The acoustic model per channel $c$ is
\begin{equation}
  y_c(t) \;=\; \sum_{i=1}^{R} \sum_{k=1}^{K}
  \operatorname{Re}\!\left[\, a_{c,i,k}(t)\, e^{\,j \phi_{i,k}(t)} \right]
  \;+\; n_c(t),
  \label{eq:model}
\end{equation}
with $\phi_{i,k}(t) = 2\pi k \int_0^t r_i(\tau)\, d\tau$, where $r_i(t)$ is
the shaft rate of rotor $i$ in rev/s, $a_{c,i,k}$ are slowly varying complex
envelopes, and $n_c$ collects broadband aerodynamic noise and any target
sources. The annotation task is to recover $\{r_i(t)\}_{i=1}^{4}$ at
sub-\revps{} accuracy over the in-flight portion of a recording.

Two corpora are used. \emph{DREGON}~\cite{strauss2018dregon} provides
8-microphone, 44.1\,kHz recordings of a quadrotor in free flight with motor
telemetry at ${\approx}929$\,Hz; the measured-speed track, where present,
serves as ground truth after standard despiking. The free-flight recordings
include speech and white-noise interferers co-recorded at low SNR.
\emph{Matrice~100} is our own single-drone corpus (DJI Matrice~100, 8-mic
ring array) whose flight-controller telemetry arrives at only
${\approx}29$\,Hz; its cruise segments provide an out-of-domain test with
ground truth adequate for scoring but not for initialization. Evaluation
pools the per-rotor mean absolute error (MAE) over the cruise portions of
held-out recordings, after permutation alignment of estimated to true rotor
tracks; capture of a rotor is declared when its trajectory error stays
within 2\revps{} of truth.

\section{Method}
\label{sec:method}

\subsection{Coupled envelope solve}

Given current trajectory estimates $r_i(t)$, form the carriers
$c_m(t) = \exp(j 2\pi \sum_{\tau \le t} k_m r_{i_m}(\tau)/f_s)$ for every
track $m = (i,k)$ and minimize over envelopes the second-generation VK
functional
\begin{equation}
  J[a] = \sum_t \Bigl| y(t) - \sum_m \operatorname{Re}[a_m(t)\,c_m(t)]
  \Bigr|^2 + \sum_m \rho_m^2 \bigl\lVert \Delta^2 a_m \bigr\rVert^2,
  \label{eq:vk2}
\end{equation}
where $\Delta^2$ is the second-difference operator and $\rho_m$ sets the
per-track selectivity, chosen from a target $-3$\,dB bandwidth by the
standard relation~\cite{tuma2000characteristics}. The normal equations of
\eqref{eq:vk2} form a block system whose diagonal blocks are banded and
whose off-diagonal blocks carry the interference terms
$\overline{c_m}\,c_n$; retaining them is what allows two tracks separated by
less than a filter bandwidth to divide energy correctly instead of both
claiming it.

Solving \eqref{eq:vk2} at audio rate for $R \times K \approx 100$ tracks is
unnecessary: the envelopes are narrowband by construction. We demodulate
each track ($z_m = \mathrm{LP}[\,y\,\overline{c_m}\,]$, zero-phase), decimate
to an envelope rate $\fsenv$ of 100--200\,Hz, and solve on the decimated
grid. The coupling coefficient survives decimation exactly for those pairs
whose instantaneous separation falls below $\fsenv/2$ --- precisely the
pairs that interfere. The decimated problem is not an approximation of
convenience: restricted to envelopes bandlimited below $\fsenv/2$ --- the
subspace the smoothness prior already enforces --- minimizing
\eqref{eq:vk2} at audio rate and minimizing its decimated counterpart are
the same problem, since every component the demodulating low-pass discards
is orthogonal to that subspace and shifts $J$ by a constant. The full-rate
signal still anchors each outer round: convergence is monitored by the
audio-rate residual ratio $\lVert y - \hat{y} \rVert^2 / \lVert y
\rVert^2$ of the full reconstruction, and carriers are re-derived at audio
rate from every updated trajectory. Tracks are partitioned into coupling groups by
union--find on this predicate; each group yields a Hermitian
positive-definite \emph{banded} system under time-major ordering of the
unknowns, factorized by banded Cholesky. Pairs whose separation never
approaches the demodulation bandwidth are pruned with no measurable effect
on trajectories. Together these choices reduce a nominal
$10^6$-unknown sparse problem to a sequence of banded solves and give the
runtime figures of Section~\ref{sec:runtime}.

\subsection{Fisher-weighted frequency refinement}

The envelope of a track demodulated with an erroneous trajectory rotates at
the residual frequency error: if $r_{i}$ is wrong by $\delta(t)$ rev/s, the
phase of $z_{i,k}$ advances by $2\pi k\,\delta(t)/\fsenv$ per envelope
sample. We estimate per-track phase slopes
$\hat{\delta}_{c,i,k}(t) = \angle\!\left( z(t{+}1)\,\overline{z(t)} \right)
\cdot \fsenv / (2\pi k)$
and fuse them across harmonics and channels into a single correction per
rotor by the weighted, smoothness-regularized problem
\begin{equation}
  \min_{\delta} \; \sum_{c,k,t} w_{c,k}(t)
  \bigl( \hat{\delta}_{c,k}(t) - \delta(t) \bigr)^2
  + \lambda \bigl\lVert \Delta^2 \delta \bigr\rVert^2,
  \qquad w_{c,k} = k^2 |z_{c,k}|^2,
  \label{eq:ifupdate}
\end{equation}
a pentadiagonal solve per rotor. The weights are the Fisher information of a
phase observation at harmonic $k$: high harmonics amplify a trajectory error
$k$-fold and are correspondingly trusted more, in proportion to their local
SNR. Each outer iteration clips the applied correction (default
$0.5\revps$) and re-derives carriers. The same envelope-phase observation
underlies the IF refinement of IAVKF~\cite{jiang2024iavkf}; the fusion
across a comb tied to one latent fundamental --- rather than one free IF per
component --- is what the rotor constraint adds, and it is where most of the
method's noise robustness originates.

\subsection{Capture annealing and blind seeding}

Phase-slope refinement is a local method: harmonic $k$ contributes usefully
only while $k\,\delta_0 < \fsenv/2$, beyond which its observation aliases.
We therefore anneal: early iterations use a small harmonic range
$K_{\max}$ and a stiff trajectory prior, and later iterations widen both as
the track locks. This is the same capture-versus-resolution compromise known
from extended Kalman comb trackers~\cite{streit1992multiharmonic}, resolved
schedule-wise rather than by inflating process noise.

Blind operation requires seeds. We scan a whitened average spectrum with a
comb matched filter over a dense fundamental grid, capping the dictionary at
1.2\,kHz: above that band, comb aliases of \emph{lower} fundamentals begin
to outscore true combs, and the cap alone repaired the corpus where naive
scanning failed outright. Because a strong rotor can shadow a weaker one at
a nearby fundamental, seeding proceeds in rounds: after each acceptance the
seeded comb is subtracted in the whitened domain and the residual re-scanned.
Finally, the multi-stage refinement ladder is wrapped in a per-track guard:
a stage whose output drops a track's comb occupancy --- the fraction of its
harmonics found above the local noise floor --- is reverted for that track.
The guard addresses a failure mode in which a re-captured track lands on a
\emph{louder, wrong} comb; raw residual confidence increases during exactly
this failure, so occupancy rather than confidence is the reliable signal.

\section{Experiments}
\label{sec:experiments}

\subsection{Telemetry-referenced accuracy}

With telemetry initialization, refinement on DREGON free-flight cruise
attains 0.73\revps{} pooled MAE against the measured-speed track --- a
figure limited in part by the reference itself, which is a command-side
signal at finite rate. This configuration is the annotation workhorse: it
runs at 0.04 real-time and its output is systematically smoother than the
telemetry it starts from.

\subsection{Blind tracking}

Table~\ref{tab:blind} summarizes the fully blind results. On the three
DREGON free-flight recordings with usable ground truth the blind tracker
reaches 0.680, 0.701 and 0.744\revps{} pooled MAE, resolving the
twin-rotor ambiguity (two rotors within one filter bandwidth) on two of the
three. On the Matrice~100 cruise segment --- a different airframe, different
microphone geometry, different telemetry pipeline, and no parameter
retuning --- it attains 1.027\revps{} pooled, per-rotor
0.67/1.19/1.22/1.03, with all four rotors captured.

\begin{table}[t]
  \centering
  \caption{Blind per-rotor tracking, pooled MAE (rev/s), cruise segments.}
  \label{tab:blind}
  \begin{tabular}{lccc}
    \toprule
    Recording & Blind MAE & Rotors captured & Twins resolved \\
    \midrule
    DREGON ff-1 & 0.680 & 4/4 & yes \\
    DREGON ff-2 & 0.701 & 4/4 & yes \\
    DREGON ff-3 & 0.744 & 4/4 & no \\
    Matrice 100 & 1.027 & 4/4 & --- \\
    \bottomrule
  \end{tabular}
\end{table}

Three components are load-bearing, and each is isolatable. Removing the
scan-band cap returns the Matrice recording to a 22\revps{} failure (the
scan seeds an alias); removing the residual re-scan leaves its fourth rotor
permanently unseeded (3.24\revps{} pooled --- the guard cannot invent a
missing seed); disabling the stage guard lets late refinement re-capture one
rotor onto a louder neighboring comb (9.7\revps{} on that rotor) while
every internal confidence measure improves.

\subsection{Comparison with learned regressors}

We compare against the strongest neural per-rotor regressors we could train:
convolutional--recurrent and convolutional--transformer architectures over
log-magnitude and instantaneous-frequency spectrogram inputs, trained on
online-mixed real recordings with matched augmentation, and multipitch
salience baselines adapted from music transcription
\cite{bittner2022basicpitch}. On the identical evaluation clips the best
neural model reaches 2.5\revps{} on DREGON cruise and 1.4--2.3\revps{}
cross-drone --- a factor 2--4 behind the blind tracker, despite the tracker
using no training data at all. The gap is consistent with our activation
analyses of those models, which show them anchored to amplitude cues rather
than comb geometry; phase-coherent tracking is precisely what they do not
learn. We regard closing this gap with causal learned trackers as open
work, and the present method as, among other things, the annotation source
such work requires.

\subsection{Runtime}
\label{sec:runtime}

On a single laptop CPU core the telemetry-refinement configuration runs at
real-time factor 0.037 and the blind configuration at 0.35. The dominant
costs are the banded envelope solves and full-length demodulation FFTs; the
time-major banded factorization is 5--10$\times$ faster than a general
sparse solve of the same systems, and far-pair pruning removes the majority
of demodulation work in richly coupled configurations at no measurable
accuracy cost. Batch annotation of an hour-long corpus is thus a matter of
minutes.

\section{Conclusion}

Order tracking, developed for shafts with tachometers, transfers to blind
acoustic annotation of multirotor flight --- provided the solver couples the
interlaced combs it will actually meet, fuses frequency evidence across each
comb rather than per line, and is seeded and guarded by procedures that
respect how comb scans fail. The resulting tracker annotates public corpora
without telemetry at sub-\revps{} accuracy and outpaces trained neural
regressors on their own evaluation data by a wide margin.

Limitations delimit the claim. The method is non-causal and, at 0.35
real-time blind, an offline tool; the twin ambiguity survives on one of
three DREGON recordings; the guard can only protect seeds that exist, so a
comb entirely absent from the scan band stays untracked; and cruise flight,
with its slowly varying speeds, is the favorable regime --- aggressive
maneuvering shrinks every capture margin involved. Within those limits, the
method turns rotor-speed side information from a dataset rarity into
something that can be manufactured for any recording of a flying machine.

\bibliographystyle{IEEEtran}
\bibliography{references}

\end{document}
